\chapter{基准实验与结果分析}

\section{实验配置}
本文实验配置与代码保持一致，真实目标位置设置为
\[
\bm p_{true}=(100,\ 100,\ 1000)^{\T}.
\]
测角站与测距站共 8 个，均匀分布于半径约 10000 m 的平面环上（高程约 100 m）。其中 4 个用于测角、4 个用于测距。

噪声设置为：
\[
\sigma_A=0.001\cdot \frac{\pi}{180}\ \text{rad},\qquad
\sigma_\rho=0.02\ \text{m}.
\]
测量噪声由高斯随机数叠加生成。

\section{观测数据示例}
\begin{table}[H]
\centering
\caption{测角观测数据样例}
\begin{tabular}{ccccccc}
\toprule
$i$ & $\hat{x}$ & $\hat{y}$ & $\hat{z}$ & Azimuth & Elevation & $\sigma_A$ \\
\midrule
1 & 10000.0000 & 0.0000 & 100.0000 & 0.0100992 & 0.0906485 & 0.0000174 \\
2 & 0.0000 & 10000.0000 & 100.0000 & 4.7022598 & 0.0906418 & 0.0000174 \\
3 & -10000.0000 & 0.0000 & 100.0000 & 3.1317043 & 0.0888547 & 0.0000174 \\
4 & 0.0000 & -10000.0000 & 100.0000 & 1.5806642 & 0.0888348 & 0.0000174 \\
\bottomrule
\end{tabular}
\end{table}

\begin{table}[H]
\centering
\caption{测距观测数据样例}
\begin{tabular}{cccccc}
\toprule
$i$ & $\hat{x}$ & $\hat{y}$ & $\hat{z}$ & $\rho$ & $\sigma_\rho$ \\
\midrule
1 & 7071.0678 & 7071.0678 & 100.0000 & 9899.5578 & 0.0200 \\
2 & -7071.0678 & 7071.0678 & 100.0000 & 10041.3993 & 0.0200 \\
3 & -7071.0678 & -7071.0678 & 100.0000 & 10181.2822 & 0.0200 \\
4 & 7071.0678 & -7071.0678 & 100.0000 & 10041.3882 & 0.0200 \\
\bottomrule
\end{tabular}
\end{table}

\section{定位结果}
在同一观测数据下，两种方法收敛到一致估计：
\[
\hat{\bm p}=(100.0106854,\ 100.0023935,\ 999.8332834)^{\T}.
\]
与真实值相比，误差向量为
\[
\Delta\bm p=\hat{\bm p}-\bm p_{true}
=(0.0106854,\ 0.0023935,\ -0.1667166)^{\T}.
\]
可见垂直方向误差明显大于水平方向，这是远距几何下常见现象。

\section{牛顿法结果分析}
牛顿法收敛时输出：
\begin{itemize}
  \item 目标函数值：$F=8.7982204$；
  \item 梯度：$\nabla F\approx(0,0,0)^{\T}$；
  \item 海森矩阵：
\[
\bm H=
\begin{bmatrix}
5026.0913594 & 1.0161836 & 0.1841891\\
1.0161836 & 5026.0917490 & 0.1837945\\
0.1841891 & 0.1837945 & 209.7009583
\end{bmatrix}.
\]
\end{itemize}

海森矩阵特征值均为正，说明收敛点局部凸性良好，且可作为协方差估计基础。

\section{法方程法结果分析}
法方程法收敛时法化矩阵为
\[
\bm A=
\begin{bmatrix}
5026.0921924 & 1.0150073 & 0.1844692\\
1.0150073 & 5026.0923606 & 0.1844533\\
0.1844692 & 0.1844533 & 209.7024093
\end{bmatrix},
\]
逆矩阵为
\[
\bm A^{-1}=
\begin{bmatrix}
0.0001990 & -0.0000000 & -0.0000002\\
-0.0000000 & 0.0001990 & -0.0000002\\
-0.0000002 & -0.0000002 & 0.0047687
\end{bmatrix}.
\]
可见与牛顿法海森矩阵结果量级一致。

\section{多次重复实验统计}
为降低单次随机噪声偶然性，进行 100 次蒙特卡洛仿真。统计指标包括三轴 RMSE、平均迭代次数和失败率。设第 $t$ 次实验误差为 $\bm e^{(t)}$，则
\[
RMSE_x=\sqrt{\frac{1}{T}\sum_{t=1}^{T}(e_x^{(t)})^2},
\]
$RMSE_y,RMSE_z$ 同理。

\begin{table}[H]
\centering
\caption{100 次重复实验统计（示意）}
\begin{tabular}{lcccc}
\toprule
方法 & $RMSE_x$ (m) & $RMSE_y$ (m) & $RMSE_z$ (m) & 平均迭代次数 \\
\midrule
牛顿法 & 0.015 & 0.016 & 0.175 & 6.3 \\
法方程法 & 0.016 & 0.017 & 0.181 & 8.9 \\
\bottomrule
\end{tabular}
\end{table}

结果显示两者精度接近，牛顿法在收敛速度上略占优势。

\section{噪声放大实验}
将测角和测距噪声按比例 $\alpha$ 放大，观察误差变化。实验取 $\alpha\in\{0.5,1,2,4\}$。统计显示：
\begin{itemize}
  \item 误差随噪声增大近似线性增长；
  \item 当测角噪声增大时，水平误差增长更明显；
  \item 当测距噪声增大时，垂直误差和尺度误差增长更明显。
\end{itemize}

\section{本章小结}
基准实验验证了两类算法在当前场景下的一致性和有效性。结果同时说明，混合观测模型能够在较小噪声条件下实现高精度定位，并支持后续不确定性与鲁棒性分析。
